% Nội dung chân trang dùng chung cho các kiểu trang.
\makeatletter
\newcommand{\pageauthor}{{\color{base800}\small
        \raisebox{0.2ex}\iconUserTie\, \@author}}
\makeatother

% Hiển thị số trang trong một dải màu.
\newcommand{\pagenumber}{%
    \begin{tikzpicture}[
            baseline=(pagenum.base),
            overlay
        ]
        \node[
            anchor=east,
            left color=white,
            right color=base200,
            inner xsep=16pt,
            inner ysep=8pt
        ] (pagenum) at (16pt,0) {\color{base800}\fontseries{b}\selectfont\thepage};
        \path[draw=none,fill=base600]
        (pagenum.south east)
        -- ++(-4pt,4pt)
        -- ($(pagenum.north east) + (-4pt,-4pt)$)
        -- (pagenum.north east)
        -- cycle;
    \end{tikzpicture}%
}

\newcommand{\pagefooter}{
    \fancyfoot[L]{\pageauthor}
    \fancyfoot[R]{\pagenumber}
}

% Định dạng đầu và chân trang mặc định.
\pagestyle{fancy}
\fancyhf{}

% Đường kẻ ngăn cách đầu trang với nội dung.
\renewcommand{\headrulewidth}{0.1pt}
\renewcommand{\headrule}{
    \color{base700}
    \hrule width\headwidth height\headrulewidth
    \vskip-\headrulewidth
}

% Tên mục hiện tại ở trái; số điện thoại ở phải.
\fancyhead[L]{\color{base800} \leftmark}
\fancyhead[R]{\color{base800} {\small\iconMobile}\, \phonenumber}
\pagefooter

% Trang đầu bài học bỏ phần đầu trang nhưng vẫn giữ chân trang.
\fancypagestyle{sectionfirst}{
    \fancyhf{}
    \renewcommand{\headrulewidth}{0pt}
    \pagefooter
}
